
\begin{proposition}[Actual minimalization at the largest canonical order]\label{cdp:st:C.2.UN.77}
Let \(f:X\to\Delta\) be a proper degeneration of connected
complex two-tori with smooth total space and full reduced
simple-NC central fibre.  Normalize its canonical generator
as in \ref{cdp:st:C.2.UN.75} and put
\[
        \operatorname{div}(\Omega)=\sum k_iV_i,\qquad
        \min k_i=0,\qquad M=\max k_i .
\tag{C.HM.1}\label{cdp:eq:C.HM.1}
\]
There is a finite overlap-compatible modification record,
over the same disc and unchanged over its punctured part,
whose endpoint has smooth total space and full reduced
normal-crossing fibre, such that every component with
canonical order \(M\) and Hopf minimal model is itself a
smooth minimal Hopf surface.  These Hopf components are
embedded, their boundary curves are disjoint smooth
elliptic curves, and no triple point meets them.
Other endpoint components may retain the two normalized
branches of an NC self-intersection.

All original Hopf-birational prime divisors survive, with
their original orders.  Every whole-component removal block
removes a ruled component of positive order \(q\);
its preliminary additions, when needed, are finitely many
ruled components of order \(q+1\le M\).  There are no such
additions when \(q=M\).  Thus the retained reduced degeneration stages have
no global component of order greater than \(M\).
Auxiliary common blowups in a type-I or type-II diagram
retain their actual multiplicities and orders; the bound
just stated does not apply to those auxiliary divisors.

The record retains its actual centres, local contractions,
normalization branches and restored overlap maps.  Its
reverse recovers the original family.  For a given compact
family the same operations take place on that compact
base, by the actual annular comparison.  Their Euler
changes are exactly the normal-cone terms of \ref{cdp:st:C.2.UN.75}.

The endpoint assertion concerns the maximal-order Hopf
components.  It makes no holomorphic inverse-logarithm
assertion for an arbitrary neighbourhood of a remaining
Hopf leaf.
\end{proposition}

\begin{proof}
\par\medskip\noindent\textbf{The surface divisor and coefficient-one curves.}
At every smooth-total NC stage, let \(H\) be the smooth
normalization of a component of order \(M\) whose minimal
model is Hopf.  The actual adjunction map, with both
preimages of each self-double curve retained, gives
\[
       -K_H=D_H=\sum_C a_C C,\qquad
       a_C=M-k_{\mathrm{other}(C)}+1\ge1 .
\tag{C.HM.2}\label{cdp:eq:C.HM.2}
\]
A self-double branch has coefficient one.  The equality
is the restriction of the common meromorphic volume
form; it retains its line and divisor, rather than only
its numerical class.

Let \(\pi:H\to H_{\min}\) be the actual sequence of point
blowdowns to its Hopf minimal surface.  A map from
\(\PP^1\) to a Hopf surface is constant: lift it to
\(\C^2\setminus\{0\}\) and use its two constant coordinate
functions.  Thus every rational curve on \(H\) is one
of the finitely many exceptional curves of \(\pi\).
Every compact curve on \(H_{\min}\) is a smooth elliptic
curve.  Indeed rational normalization is impossible,
and \(H^2(H_{\min},\Q)=0\) and adjunction give arithmetic
genus one; the normalization genus therefore equals
one and its delta invariant is zero.  Distinct such
curves are disjoint, since their intersection number
is zero and all local intersection multiplicities are
positive.

Push forward \eqref{cdp:eq:C.HM.2} through the point blowdowns.  At
each reverse point blowup \(\sigma\) one has the exact
divisor formula
\[
             D_{\mathrm{new}}=\sigma^*D_{\mathrm{old}}-E.
\tag{C.HM.3}\label{cdp:eq:C.HM.3}
\]
Every intermediate divisor is effective: a negative
coefficient would persist on its strict transform in
the final effective divisor.  Starting from the
disjoint elliptic support on \(H_{\min}\), all these
blowups consequently have centres on the positive
support.  The positive support stays SNC.  At a smooth
point of a component of coefficient \(a\), the new
coefficient is \(a-1\); at a node of coefficients
\(a,b\), it is \(a+b-1\).  Coefficient-zero exceptional
curves are omitted from the positive support, but are
retained as curves of the surface.

It follows by this actual point-map induction that
\[
 \begin{gathered}
 \text{two coefficient-one positive curves never meet;}\\
 \text{a rational coefficient-one curve is a leaf
 of the positive boundary,}\\
 \text{and an elliptic coefficient-one curve is
 disjoint from that boundary.}
 \end{gathered}
\tag{C.HM.4}\label{cdp:eq:C.HM.4}
\]
For the first assertion a smooth blowup creates
\((a,a-1)\), never \((1,1)\); a node blowup could
create a new \((1,1)\) only from an old \((1,1)\).
For the second assertion a new rational coefficient
one is first created at a smooth point of coefficient
two, and has valence one.  Subsequent node blowups
preserve that valence, and a smooth blowup at its
points has coefficient zero.  An original coefficient
one elliptic curve likewise gains no positive boundary
at a smooth blowup.

These facts also rule out a rational self-double curve
on a maximal Hopf component at any NC stage in the
retained class.  Its two normalized preimages, as
specified by \cite[\S4.11.2(III)]{Kulikov1977}, have
coefficient one.  Each is rational and has exactly one
boundary intersection by \eqref{cdp:eq:C.HM.4}; hence the image has
one triple point.  Each rational curve has negative
self-intersection on the blown-up Hopf, whereas the
actual two-normal-bundle formula would give
\[
              C_1^2+C_2^2=-1,
\tag{C.HM.5}\label{cdp:eq:C.HM.5}
\]
an impossibility.  A triple point with three branches
of this same component would already give two
coefficient-one curves meeting, also impossible.
The modifications below can create new self-double
curves only from rational exceptional curves.
They therefore create no elliptic self-double curve
on a maximal Hopf starting from the given embedded
components.  At every subsequent smooth-total NC
stage these Hopf components are consequently embedded.

\par\medskip\noindent\textbf{The exact last-point alternatives.}
Suppose such an \(H\) is not minimal and choose a last
point blowup in the actual map \(\pi\).  Its exceptional
curve \(E\) is a smooth \((-1)\)-curve.  Formula
\eqref{cdp:eq:C.HM.3} leaves exactly three cases.

First, \(E\) is not a boundary curve.  It was created
at a smooth point of a coefficient-one boundary \(C\),
and meets that curve once.  The ambient normal
sequence is
\[
 0\longrightarrow\OO_{\PP^1}(-1)
 \longrightarrow N_{E/X}
 \longrightarrow\OO_{\PP^1}(-1)\longrightarrow0 .
\tag{C.HM.6}\label{cdp:eq:C.HM.6}
\]
Its splitting and the ordinary two-ruling diagram of
\ref{cdp:st:C.2.UN.75} give the actual type-I move.  The neighbour across
\(C\) has order \(M\).  It is not Hopf-birational.
If \(C\) is elliptic, its square on \(H\) is at most
\(-1\) and it contains no triple point; its square on
the neighbour is therefore positive.  If \(C\) is
rational, it was already exceptional before this last
blowup, so its current square is at most \(-2\);
\eqref{cdp:eq:C.HM.4} gives one triple point, and its square on
the neighbour is again positive.  Such squares are
impossible for the corresponding elliptic or rational
curves on a blown-up Hopf.  The move removes one
point blowup from \(H\), adding it to a non-Hopf
component.

Second, \(E\) was created at a node of two positive
boundary curves of coefficients \(a,b\).  It has
coefficient \(a+b-1\) and two triple points.  The
two-normal-bundle equation makes its square on the
other component also \(-1\), so the actual type-II
move of \ref{cdp:st:C.2.UN.75} applies.  It blows down \(E\) on both
sides and blows up the two opposite normalized
branches.  If an opposite component has order \(M\),
the associated curve \(C\) on \(H\) has coefficient
one.  It is a rational leaf by \eqref{cdp:eq:C.HM.4}, and its
square after the last node blowup is at most \(-2\).
Its square on that opposite component is
\[
                 -1-C_H^2\ge1 .
\tag{C.HM.7}\label{cdp:eq:C.HM.7}
\]
Hence that opposite component is not Hopf-birational.
This proves the same assertion when both opposite
branches belong to one global component: their actual
normalization branches are used in \eqref{cdp:eq:C.HM.7}.

Let
\[
       h=\sum_{\substack{k_i=M\\
                     V_i\ \mathrm{Hopf\text{-}birational}}}
                              b_2(V_i).
\tag{C.HM.8}\label{cdp:eq:C.HM.8}
\]
At the smooth NC stages just described, this is the
total number of surface point blowups on these Hopfs.
Each of the two moves above strictly decreases \(h\).
No assertion that \(h\) decreases inside the longer
ruled-removal blocks below is needed.

Third, \(E\) was created at a smooth point of a
positive curve of coefficient \(a\ge2\).  It has
coefficient \(a-1\) and one triple point.  Write \(Q\)
for its neighbour and \(P\) for the third component.
The exact normal degrees and canonical orders are
\[
 \begin{split}
 E_H^2&=-1,\qquad E_Q^2=0,\\
 q:=k_Q&=M-a+2,\qquad k_P=q-1,\\
                         &1\le q\le M .
 \end{split}
\tag{C.HM.9}\label{cdp:eq:C.HM.9}
\]
The complete rational curve of square zero gives
the actual ruling on the compact normalization of
\(Q\).  For clarity, its local complete deformations
are disjoint, since their intersection number is zero;
the resulting compact curve family gives the ruling
after the point blowdowns of its singular fibres.
Equivalently this is the ruled-surface comparison in
\cite[\S4.5]{Kulikov1977}.  Its one boundary intersection
is the section belonging to \(P\).  Any other
horizontal boundary would meet a general complete
fibre in this same local family, which it does not.
All other boundary curves on \(Q\) are vertical.
This is the precise ruled section-and-fibre
configuration, including any nonminimal fibres.

\par\medskip\noindent\textbf{Ruled removal and cleanup of its singular charts.}
We use the actual finite maps in
\cite[\S\S4.10.7--4.10.9,4.14]{Kulikov1977}, with
the missing-fibre correction in
\cite[\S8]{Kulikov1981}.  Their overlaps and proper
composite graphs are those verified in \ref{cdp:st:C.2.UN.75}.
They remove the ruled component in \eqref{cdp:eq:C.HM.9}.
If its boundary is disconnected, they first insert
the specified missing fibre components; these
exceptional ruled surfaces have order \(q+1\).
When \(q=M\), the canonical divisor on its
normalization has only negative boundary
coefficients, and \cite[Lemma2.19]{Kulikov1977}
makes that boundary connected.  No preliminary
addition is then made.  Thus all added orders are
at most the original fixed \(M\).

The result of this removal need not immediately
have smooth total space.  We retain precisely the
finite sigma-atlas of
\cite[\S4.11.2]{Kulikov1977}: its singular sets are
contained in finitely many disjoint charts; in each
chart the recorded local exceptional branches
contract successively to a smooth NC chart.  We
now describe an actual cleanup, before applying
another last-point argument to a Hopf.

If one of these lists is nonempty, choose a
component of largest canonical order among the
globally occurring last locally blown-up branches.
Use an auxiliary resolution \(a:\widehat X\to X'\)
as in \cite[\S4.13]{Kulikov1977} to read its
old strict component \(\widehat V\).  This
resolution is a diagnostic model for the
hypotheses of Lemma4.14; it is not inserted
as an additional degeneration stage in the
modification record.  In particular its new
exceptional divisors receive no card labels,
and no bound on their canonical orders is used.
Each branch has a complete local exceptional
\(\PP^1\)-fibre.  These complete fibres extend to
a ruling of the old smooth surface
\(\widehat V\).  In the
generic local exceptional chart there is exactly
one horizontal boundary section.  Completeness
of those fibres shows that there can be no
additional global horizontal boundary; all the
other boundary, including the diagnostic
exceptional curves over the finite exceptional
points, is vertical.  This supplies
conditions (a), (b) and (d) of
\cite[Lemma4.14]{Kulikov1977}, by their actual
normalization and section maps.

If another curve on this chosen component had
been locally blown up, its exceptional branch
would have order one larger.  If that child
is not last, follow its finite local descendant
list.  The order increases at each step, so
a last descendant has strictly larger order
than the chosen component.  This contradicts
the choice of largest last-branch order.
Thus the choice gives condition
(e).  This is the exact comparison in
\cite[\S4.16.2]{Kulikov1977}.  The reordering in
its \S4.12 makes its specified branches last
without changing their restored overlaps.
The correction in \cite[\S8]{Kulikov1981}
inserts any missing vertical links in the
specified smooth open charts, leaving their
overlap fixed.  Lemma4.14 then removes this
global ruled component, and the result again
has the same kind of sigma-atlas.

Every cleanup removal has positive order:
a locally exceptional component has the
actual order of its parent plus one.  It
adds only finitely many components at the
next order, and none at \(M\).  Here we
use precisely the maximal-old-order clause
of \cite[\S8]{Kulikov1981}, including its
last-locally-blown-up case: its boundary
is connected and no missing links are
inserted.  This does not assert that
every divisor on the auxiliary resolution
has order at most \(M\).  A Hopf
component cannot be one of these locally
exceptional global components: its complete
local rational fibres would give rational
curves through an open set of its compact
normalization, contrary to the finite
exceptional-curve description above.
The strict Hopf divisors survive by the
proper graph argument in \ref{cdp:st:C.2.UN.75}.

This cleanup terminates without inserting an
ambient resolution into the record or estimating
orders of
arbitrary resolution divisors.  Label every
current global component by its order in
\(\{0,\ldots,M\}\).  A removal deletes a
positive label \(q\), adds finitely many
labels \(q+1\) when \(q<M\), and adds none
when \(q=M\).  There are finitely many
removals at level one, hence finitely many
objects ever created at level two, and
inductively at every level through \(M\).
This is the finite game of
\cite[\S9]{Kulikov1981}, now applied also
to the cleanup removals.

When cleanup stops, no recorded locally
exceptional branch remains.  The defining
contraction lists of the sigma-atlas are
empty; their terminal smooth NC charts
are therefore the current charts.
The total space is smooth and its full
fibre is reduced NC.  Its normalization
components are smooth.  Other global
components may retain self-incidences.
At a maximal Hopf the intrinsic argument
\eqref{cdp:eq:C.HM.2}--\eqref{cdp:eq:C.HM.5} excludes rational
self-incidences.  No elliptic self-incidence
was introduced: type-I and type-II changes
use rational exceptional curves; in a
local ruled contraction the horizontal
parent has order one lower, so a
maximal-order Hopf can only meet that
contraction in vertical rational curves.
Thus the maximal Hopfs are again smooth
embedded components, and their ordinary
point maps to their Hopf minimal models
are available for the next step.

\par\medskip\noindent\textbf{Termination of the whole procedure.}
Perform a type-I or type-II step whenever
the chosen last point gives the first or
second case.  In the third case perform
the finite ruled removal and cleanup
just specified.  Across the entire
procedure there are only finitely many
divisorial removal blocks: the same
fixed-level game counts all their
positive component labels together.
Between these blocks every step strictly
decreases the finite integer \eqref{cdp:eq:C.HM.8}.
More explicitly, if \(c_j\) counts the
retained global components of order \(j\),
the block-level potential is
\[
                 (c_1,c_2,\ldots,c_M,h)
                 \quad\hbox{in lexicographic order.}
\tag{C.HM.10}\label{cdp:eq:C.HM.10}
\]
A whole-component removal decreases its
first changed \(c_j\), whatever finite
next-level additions or Hopf point
changes occur inside the block.
A type-I or type-II step chosen on a
last Hopf point leaves all the \(c_j\)
unchanged and decreases \(h\).
Lexicographic order on this fixed
finite product of nonnegative integers
has no infinite strictly descending
sequence, by the same successive-coordinate
induction.  This potential is attached
to the specified finite blocks, not to
every auxiliary arrow in their graphs.
An infinite procedure would therefore
have either infinitely many divisorial
removals, contradicting that game, or
an infinite tail of strict decreases of
\(h\).  Both are impossible.

At the endpoint a nonminimal maximal
Hopf would still have a last point
blowup and one of the three alternatives,
so every such Hopf is minimal.  Its
curves are the disjoint smooth elliptic
curves proved above, and hence no
triple point meets it.  The original
Hopf primes survive with unchanged
orders by \ref{cdp:st:C.2.UN.75}, and each additional
prime used in the record is ruled.
The same actual proper graphs and
restored overlap isomorphisms give
the asserted reversible comparison.
Their annular gluing and normal-cone
Euler terms are exactly those of \ref{cdp:st:C.2.UN.75}.
This proves the statement in its
specified scope.
\end{proof}


